\documentclass[a4paper]{article}

\usepackage[fontset=fandol]{ctex}
\usepackage{amsmath, amssymb}
\usepackage{graphicx}
\usepackage[margin=2.45cm]{geometry}
\usepackage[most]{tcolorbox}
\usepackage{hyperref}
\usepackage{booktabs}
\usepackage{float}
\usepackage{array}
\usepackage{xcolor}
\usepackage{enumitem}
\usepackage{caption}

\setlength{\parindent}{2em}
\setlength{\parskip}{0.35em}
\setlength{\emergencystretch}{3em}
\linespread{1.06}
\sloppy
\raggedbottom
\vfuzz=4pt

\newcolumntype{L}[1]{>{\raggedright\arraybackslash}p{#1}}
\newcolumntype{C}[1]{>{\centering\arraybackslash}p{#1}}

\newcommand{\Dtrain}{\mathcal{D}^{\mathrm{train}}}
\newcommand{\Dquery}{\mathcal{D}^{\mathrm{query}}}
\newcommand{\Task}{\mathcal{T}}
\newcommand{\flearner}{F_{\phi}}
\newcommand{\loss}{L}
\newcommand{\innerloss}{\ell}
\newcommand{\grad}{\nabla}

\hypersetup{
    colorlinks=true,
    linkcolor=blue!60!black,
    urlcolor=blue!60!black
}

\setlist[itemize]{leftmargin=2.2em,itemsep=0.22em,topsep=0.25em}
\setlist[enumerate]{leftmargin=2.4em,itemsep=0.22em,topsep=0.25em}
\setlist[description]{leftmargin=3.1em,itemsep=0.18em,topsep=0.25em}

\newtcolorbox{knowledgebox}[1]{
    enhanced, colback=blue!5!white, colframe=blue!75!black, colbacktitle=blue!75!black,
    coltitle=white, fonttitle=\bfseries, title={#1},
    attach boxed title to top left={yshift=-2mm, xshift=2mm},
    boxrule=1pt, sharp corners
}

\newtcolorbox{importantbox}[1]{
    enhanced, colback=yellow!10!white, colframe=yellow!80!black, colbacktitle=yellow!80!black,
    coltitle=black, fonttitle=\bfseries, title={#1}, sharp corners
}

\newtcolorbox{warningbox}[1]{
    enhanced, colback=red!5!white, colframe=red!75!black, colbacktitle=red!75!black,
    coltitle=white, fonttitle=\bfseries, title={#1}, sharp corners
}

\newcommand{\notetitle}{元学习（二）：万物皆可 Meta}
\newcommand{\notesubtitle}{从初始化、优化器、架构到数据处理与 few-shot 应用}
\newcommand{\noteauthors}{基于李宏毅老师《机器学习 2025》课程内容整理}
\newcommand{\notedate}{\today}
\newcommand{\videochannel}{李宏毅深度学习课堂（Bilibili）}
\newcommand{\videoduration}{约 31 分 36 秒}
\newcommand{\videourl}{https://www.bilibili.com/video/BV1TAtwzTE1S?p=56}

\newcommand{\framefig}[4]{%
\begin{figure}[H]
    \centering
    \includegraphics[width=#1\textwidth]{#2}
    \caption{#3\protect\footnotemark}
\end{figure}
\footnotetext{视频画面时间区间：#4。}
}

\begin{document}
\pagestyle{empty}

\begin{center}
    \vspace*{1.0cm}
    {\Huge\bfseries \notetitle\par}
    \vspace{0.35cm}
    {\LARGE \notesubtitle\par}
    \vspace{0.75cm}
    \includegraphics[width=0.62\textwidth]{video.jpg}\par
    \vspace{0.75cm}
    {\large \noteauthors\par}
    {\large \notedate\par}
    \vspace{0.75cm}
    \begin{tcolorbox}[width=0.86\textwidth,center,sharp corners,
                      colback=black!3!white,colframe=black!50!white]
        \centering
        \textbf{视频信息}\\[3pt]
        频道：\videochannel \\
        时长：\videoduration \\
        链接：\href{\videourl}{\videourl}
    \end{tcolorbox}
\end{center}

\newpage
\tableofcontents
\newpage
\pagestyle{plain}

% =====================================================================
\section{这一讲的问题：学习算法里什么可以学}
% =====================================================================

P55 建立了元学习的基本框架：普通机器学习学习函数 $f_\theta$，元学习学习学习算法 $F_\phi$。P56 进一步追问：一个学习算法里面到底哪些部分可以被放进 $\phi$？课程标题“万物皆可 Meta”不是说什么都随便套元学习，而是说许多原本靠人工设定的训练环节，都可以被写成可学习对象。

\framefig{0.84}{figures/fig_01_learnable_question.jpg}
{P56 的核心问题：What is learnable in a learning algorithm? 也就是在一个学习算法中，哪些组件可以由跨任务经验学习出来。}
{00:00:00--00:00:15}

先把普通梯度下降写成一个流程。给定网络结构、初始参数、训练资料和更新规则，算法反复计算梯度并更新参数：
\[
    \theta^{t+1}
    =
    \theta^{t}
    -
    \eta \nabla_{\theta}\ell(\theta^{t};\Dtrain).
\]
在这个式子里，很多东西都不是数据自动决定的：网络结构要人选，初始参数通常随机抽样，学习率 $\eta$ 要调，优化器形式要选，训练资料还要经过 augmentation 或 weighting。

\framefig{0.84}{figures/fig_02_gradient_descent_components.jpg}
{梯度下降流程中可以被元学习改造的组件：network structure、initialization、update rule、gradient computation 与 training data 处理都可能成为 $\phi$ 的一部分。}
{00:01:15--00:01:30}

因此，本讲可以看成一张“可 Meta 化对象”的地图：
\begin{itemize}
    \item 学初始参数：learning to initialize，代表方法是 MAML。
    \item 学优化器：learning to optimize，学习更新规则、学习率、动量等。
    \item 学网络架构：network architecture search，把结构选择也视为可优化对象。
    \item 学数据处理：data augmentation、sample reweighting、课程式采样策略等。
    \item 学分类过程：metric-based meta-learning，把训练和测试合成一个可学习函数。
\end{itemize}

\begin{importantbox}{“Meta” 的统一视角}
    判断一个方法是否属于元学习，不要只看它是否叫 MAML 或 few-shot。更稳的判断是：它是否利用许多 tasks 来学习某种“如何学习”的规则，并希望这条规则迁移到新 task 上。
\end{importantbox}

\subsection{本章小结}

P56 的主线是把学习算法拆开：初始化、优化器、架构、数据处理、样本权重、分类流程都可以成为 $\phi$。元学习的实质是把这些原本手工设定的组件，改成可由 task-level 经验训练出来的对象。

% =====================================================================
\section{Learning to Initialize：MAML 与 Reptile}
% =====================================================================

第一类方法学习初始参数。普通训练中，初始参数通常从固定随机分布抽样；learning to initialize 认为初始点本身可以被学习。如果初始点选得好，新任务只需少量梯度更新就能得到不错模型。

\framefig{0.84}{figures/fig_03_learning_to_initialize_maml.jpg}
{Learning to initialize 的代表方法包括 MAML 与 Reptile：它们学习一个适合快速适应的初始参数，使模型面对新任务时可以用少量样本和少数更新步骤完成适应。}
{00:02:15--00:02:30}

MAML 的全名是 Model-Agnostic Meta-Learning。Model-agnostic 表示它不强依赖某种特定网络结构，只要模型可以用梯度下降更新，就可以尝试套用。一个简化的一步 MAML 过程如下。
\[
    \theta_n'
    =
    \phi
    -
    \alpha \nabla_{\phi}
    \ell_{\mathrm{train}}^n(\phi).
\]
其中 $\phi$ 是共享初始化，$\theta_n'$ 是第 $n$ 个 task 经过一次 inner update 后的参数。然后用该 task 的 query set 评价：
\[
    L(\phi)
    =
    \sum_{n=1}^{N}
    \ell_{\mathrm{query}}^n(\theta_n').
\]
最后对 $\phi$ 做 outer update。训练时，优化器不是只看某个 task 的 training loss，而是寻找一个初始化，使它经过少量 task-specific 更新后能在 query examples 上表现好。

\begin{knowledgebox}{MAML 的直觉}
    MAML 不直接要求 $\phi$ 本身就是所有 task 的最佳参数，而是要求 $\phi$ 离许多 task 的好解都不远。它学习的是“容易被少量资料推到正确方向”的起点。
\end{knowledgebox}

训练 MAML 本身也要做优化，而且可能并不稳定。课程引用了 How to train your MAML 一类工作，说明训练这种初始化并不是简单跑一次梯度下降就好；inner loop 步数、学习率、二阶梯度近似、任务采样、归一化等都会影响稳定性。

\framefig{0.84}{figures/fig_04_training_maml_curves.jpg}
{MAML 训练可能非常敏感：不同随机种子或设置会带来不同曲线，说明“学习初始化”本身也是一个需要精心优化的 meta-training 问题。}
{00:03:30--00:03:45}

\begin{warningbox}{MAML 不是魔法初始化}
    MAML 学到的是跨 training tasks 的好起点。如果 training tasks 与未来 testing task 分布差异很大，或者 inner update 的训练资料太少、噪声太大，快速适应不一定可靠。它的成功依赖 task distribution 的相似性。
\end{warningbox}

\subsection{本章小结}

Learning to initialize 把初始参数放进 $\phi$。MAML 的目标不是直接找一个万能分类器，而是找一个能被少量新任务样本快速调整的起点。它把普通梯度下降嵌入 meta-training，因此表达力强，但训练也更敏感。

% =====================================================================
\section{MAML 与 Pre-training：名字不同，边界要看清}
% =====================================================================

课程专门比较 MAML 与 pre-training。二者都可能产出一个“好的初始化”，但训练信号不同。Pre-training 往往在大量样本或 proxy task 上先训练一个模型，再把它迁移到下游任务；MAML 则显式模拟“看 support set 后快速适应到 query set”的过程。

\framefig{0.84}{figures/fig_05_maml_vs_pretraining.jpg}
{MAML 与 pre-training 都可能给测试任务提供一个好起点；差异在于 MAML 用 task/episode 结构训练快速适应，而典型 pre-training 则先在 proxy tasks 或大规模资料上训练模型。}
{00:06:00--00:06:15}

这一区别并不总是泾渭分明。现代 self-supervised pre-training、multi-task learning、domain adaptation 和 meta-learning 之间有很多交叠。课程提醒：不要过度拘泥于术语标签，更应看方法背后的训练目标与使用方式。

如果一种方法训练时没有显式模拟新 task 的少样本适应，而只是用很多资料学一个通用表示，它更像 pre-training 或 multi-task learning；如果它训练时不断采样 task，把 support/query 的适应和评价作为目标，就更像 meta-learning。

MAML 为什么有效也有两种解释：一种认为它真的学到了 rapid learning 的好初始化；另一种认为它主要学到了可以复用的 representation，inner loop 只是在已有特征上做少量调整。ANIL 相关研究指出，即使几乎不更新内层特征，效果也可能很好，这支持 feature reuse 的解释。

\framefig{0.84}{figures/fig_06_maml_rapid_learning_feature_reuse.jpg}
{MAML 的效果可能来自两个机制：rapid learning 认为初始化便于快速移动到新 task 好解，feature reuse 则认为主要价值来自可迁移表示，inner loop 可能只需调整少量头部参数。}
{00:11:00--00:11:15}

\begin{importantbox}{看机制，不只看名字}
    当一个方法声称是 meta-learning 时，应追问它到底学习了什么：是学习初始化、表示、距离函数、优化器，还是只是在多任务资料上预训练？这些机制可能组合在一起，单靠方法名很难判断。
\end{importantbox}

\subsection{本章小结}

MAML 与 pre-training 都可能产出可迁移起点，但 MAML 的目标更强调“经过少量 task 内更新后在 query set 上表现好”。实际研究中两者边界会模糊，理解时应看训练结构、适应过程和泛化目标，而不是只看论文标签。

% =====================================================================
\section{Learning to Optimize：优化器也可以学}
% =====================================================================

第二类可学习对象是 optimizer。普通梯度下降的更新规则中有许多人为设定：学习率、动量、衰减、Adam 的动量系数、RMSProp 的平滑系数等。Learning to optimize 问的是：能不能让机器从许多优化任务中学习一个优化器？

\framefig{0.84}{figures/fig_07_learned_optimizer_phi.jpg}
{Optimizer 也可以被元学习：学习率、动量、更新公式或优化器内部状态都可以成为 $\phi$；学习目标是让这个 optimizer 在一批任务上带来更快或更稳定的下降。}
{00:13:15--00:13:30}

一个抽象写法是：
\[
    \theta^{t+1}
    =
    \theta^t
    +
    g_{\phi}
    \left(
        \nabla_{\theta}\ell(\theta^t),
        h^t
    \right),
\]
其中 $g_{\phi}$ 是可学习更新器，$h^t$ 是优化器状态，例如过去梯度、动量或隐藏状态。传统优化器的更新规则由人写死；learned optimizer 则从经验中学习“看到梯度后该怎么动”。

课程提到早期工作 Learning to learn by gradient descent by gradient descent。它们会在一些训练优化任务上学优化器，再测试能否迁移到不同网络、不同层数或不同数据集。

\framefig{0.84}{figures/fig_08_learned_optimizer_results.jpg}
{Learned optimizer 的实验曲线：优化器可在若干训练任务上学习更新策略，并尝试迁移到不同网络或数据集；关键问题是泛化能力，而不只是训练任务上的下降速度。}
{00:15:00--00:15:15}

\begin{warningbox}{学优化器的主要风险}
    学到的 optimizer 可能过拟合训练任务，只在见过的 loss landscape 上好用；也可能在长时间训练时不稳定。评价 learned optimizer 时，必须看它是否能迁移到新网络、新资料、新 loss 形状，而不是只看训练任务曲线。
\end{warningbox}

\subsection{本章小结}

Learning to optimize 把更新规则本身放进 $\phi$。它不再只问“参数应该怎么变”，而是问“看到梯度和状态以后，优化器应该怎样产生更新”。这一路线很强，但泛化与稳定性是核心难点。

% =====================================================================
\section{Network Architecture Search：结构也可以 Meta}
% =====================================================================

第三类对象是网络结构。传统上，网络架构由人设计：层数、宽度、卷积核、skip connection、attention block 等都需要选择。Network Architecture Search（NAS）把架构选择变成优化问题：找一个结构，使训练出来的网络在验证或测试任务上表现好。

在强化学习式 NAS 中，一个 agent 逐步产生架构决策，例如选择 filter size、stride、layer width、skip connection 等。这个 agent 的参数就是 $\phi$，它形成一个网络后，网络再经过 within-task training，最后用验证准确率或 loss 作为 reward 更新 agent。

\framefig{0.84}{figures/fig_09_nas_reinforcement_learning.jpg}
{强化学习式 NAS：agent 的参数 $\phi$ 决定网络架构，生成架构后训练网络并取得表现，把 $-L(\phi)$ 或准确率视为 reward 来更新 agent。}
{00:17:30--00:17:45}

把 NAS 放进元学习框架后，可以看到两层训练：
\begin{itemize}
    \item \textbf{Within-task training：}给定某个架构，训练该架构下的普通网络参数。
    \item \textbf{Across-task training：}更新产生架构的 agent，使它以后能提出更好的架构。
\end{itemize}

\framefig{0.84}{figures/fig_10_nas_across_within_task.jpg}
{NAS 中的 across-task 与 within-task：agent 产生网络结构，网络经过普通训练得到性能，性能再反馈给 agent；若训练与测试任务不同，NAS 更接近 meta-learning 设定。}
{00:19:45--00:20:00}

早期 NAS 常被质疑“训练和测试都在同一数据集上调结构”，容易像是在同一 benchmark 上过度搜索。后来许多工作开始让 training tasks 与 testing tasks 不同，评估搜索策略是否能迁移到新任务，这就更符合元学习精神。

除了强化学习式 NAS，还有可微分架构搜索 DARTS。DARTS 把离散结构选择放松成连续权重，使 $\nabla_{\phi}L(\phi)$ 可以计算，从而用梯度下降直接优化架构参数。

\framefig{0.84}{figures/fig_11_darts_differentiable_nas.jpg}
{DARTS 将架构搜索近似成可微分问题，让原本离散的 architecture choice 可以通过梯度下降优化；这相当于把 network architecture 放进 $\phi$ 并令其可微。}
{00:20:15--00:20:30}

\subsection{本章小结}

NAS 把网络结构从人工设计对象变成可学习对象。强化学习式 NAS 用 agent 产生架构并根据 reward 更新，DARTS 则把结构选择变得可微。它们说明“可学习组件”不必只限于参数或优化器，连模型形状本身也可以被 meta-optimized。

% =====================================================================
\section{Data Processing：资料处理也可以学}
% =====================================================================

除了模型和优化器，训练资料的处理方式也可以 Meta。数据增强通常靠人工经验和试错：是否翻转、裁切、旋转、改变颜色、加入噪声、混合样本，都会影响模型泛化。既然 augmentation policy 会影响训练结果，就可以把 policy 当成 $\phi$ 来学。

\framefig{0.84}{figures/fig_12_learned_data_augmentation.jpg}
{Data augmentation policy 也可被学习：系统从候选操作中组合增强策略，目标是让经过增强训练的模型在目标任务或验证集上表现更好。}
{00:20:45--00:21:00}

这类方法的典型结构是：先由 policy network 产生增强操作，再用增强后的数据训练模型，最后用模型表现更新 policy。它和 NAS 很像，只是搜索对象从 network architecture 变成 data transformation。

另一类资料处理是 sample reweighting。训练集中并非所有样本都应同等对待：边界附近样本可能更有信息量，也可能只是标签噪声；容易样本可以稳定训练，困难样本可以推动决策边界。到底给哪些样本更大权重，可以人工设定，也可以学习。

\framefig{0.84}{figures/fig_13_sample_reweighting.jpg}
{Sample reweighting：给不同样本不同权重。难样本可能值得更高权重，也可能是 noisy label 而应降低权重；权重策略本身可以由 meta-learning 学出。}
{00:21:15--00:21:30}

若把样本权重写成 $w_{\phi}(x,y)$，普通训练 loss 可改为
\[
    L(\theta;\phi)
    =
    \sum_{(x,y)\in \Dtrain}
    w_{\phi}(x,y)
    \ell(f_{\theta}(x),y).
\]
外层目标则用验证集或 query set 来更新 $\phi$。这样，模型参数 $\theta$ 学分类器，meta-parameters $\phi$ 学“哪些样本更值得相信或强调”。

\begin{importantbox}{数据处理的 Meta 化很实用}
    许多工程问题的瓶颈不在网络结构，而在训练数据质量、增强策略、噪声标签和样本不平衡。把 augmentation 或 reweighting 学出来，往往比再堆一层网络更接近真实痛点。
\end{importantbox}

\subsection{本章小结}

Data processing 也能成为 $\phi$。Data augmentation 学的是如何制造训练变化，sample reweighting 学的是哪些样本应被强调或削弱。它们把元学习从“模型内部”扩展到“训练资料流”的设计。

% =====================================================================
\section{Beyond Gradient Descent：学习与分类可以合在一起}
% =====================================================================

到目前为止，很多方法仍保留“训练阶段”和“测试阶段”的分界：先用 training data 训练出某个参数 $\theta^\ast$，再把 $\theta^\ast$ 用在 testing data 上。课程接着问：能否把这个分界拿掉，让一个函数同时读入训练资料和测试资料，直接输出测试资料的分类结果？

\framefig{0.84}{figures/fig_14_learning_to_compare.jpg}
{Learning to compare 或 metric-based meta-learning：不一定先显式训练出 $\theta^\ast$，而是让一个函数同时接收 support examples 与 query example，直接输出 query 的类别。}
{00:24:30--00:24:45}

这就是 metric-based meta-learning 的核心直觉。系统学习一个比较函数、嵌入空间或关系网络，使 query example 可以与 support examples 比较。常见路线包括 Siamese Network、Matching Networks、Prototypical Networks、Relation Networks 等。

这类方法仍然符合元学习框架，只是学习算法 $F_\phi$ 的形状不同。它不一定输出一个显式分类器参数 $\theta^\ast$，而是输出 query 的预测：
\[
    \hat{y}_{q}
    =
    F_{\phi}
    \left(
        \Dtrain,
        x_q
    \right).
\]
其中 $\Dtrain$ 是 support set，$x_q$ 是 query input。训练目标仍然来自许多 episodes：给一组 support 和 query，要求模型正确分类 query。

\begin{knowledgebox}{为什么 metric-based 方法适合 few-shot}
    当每类样本很少时，完整训练一个高维分类器可能不稳定。Metric-based 方法改成学习“如何比较”，让类别由 support examples 临时定义。新类别出现时，不必重新训练整个网络，只需把新类别 support examples 放进比较函数。
\end{knowledgebox}

\subsection{本章小结}

Beyond gradient descent 的思路是：学习算法不一定必须是“梯度下降产生参数”。Metric-based meta-learning 让模型直接学习 support-query 的比较机制，把学习和分类合并进一个函数中。

% =====================================================================
\section{应用：Few-shot Image Classification 与 Omniglot}
% =====================================================================

元学习最经典的应用场景之一是 few-shot image classification。任务设定通常写成 $N$-way $K$-shot：每个 episode 有 $N$ 个类别，每个类别只有 $K$ 个 support examples。模型要根据这些少量样本判断 query 属于哪个类别。

\framefig{0.84}{figures/fig_15_n_way_k_shot.jpg}
{Few-shot image classification 的 $N$-way $K$-shot 设定：每个 task 有 $N$ 个类别，每类只有 $K$ 个 support examples，模型需判断 query 属于哪一类。}
{00:26:00--00:26:15}

经典 benchmark 是 Omniglot。它包含 1623 个字符类别，每个字符有 20 个样本。由于类别多、每类样本少，非常适合构造大量 few-shot tasks。

\framefig{0.84}{figures/fig_16_omniglot_dataset.jpg}
{Omniglot 数据集包含 1623 个字符类别，每类约 20 个样本；在元学习里，它常被用作快速构造 few-shot classification episodes 的 benchmark。}
{00:27:15--00:27:30}

构造任务时，通常先把字符类别拆成 meta-training characters 与 meta-testing characters。然后从 meta-training characters 中随机抽取 $N$ 个类别，每类抽取 $K$ 个 support examples，再抽取 query examples，形成一个 training episode。Meta-testing 时则从完全未见过的 characters 中构造 episodes。

\framefig{0.84}{figures/fig_17_omniglot_task_sampling.jpg}
{Omniglot 上的 episode 构造：先拆分 training/testing characters，再从 training characters 中采样 $N$ 个类别、每类 $K$ 个 support examples，组成一个 meta-training task。}
{00:29:45--00:30:00}

元学习并不只适用于图像。课程最后展示了语音、自然语言处理等领域中的应用表：learning to initialize、learning to compare、NAS 等思想都可以迁移到不同输入形式与任务类型中。

\framefig{0.84}{figures/fig_18_meta_learning_applications_table.jpg}
{Meta-learning 应用不局限于 Omniglot 或图像分类；课程列出语音、NLP 等领域中 learning to initialize、learning to compare 与其他 meta-learning 方法的应用。}
{00:30:45--00:31:00}

\begin{warningbox}{Benchmark 不能代表全部}
    Omniglot 很适合教学和快速验证，但真实任务往往有类别不平衡、域偏移、噪声标签、任务边界不清等问题。元学习方法能否落地，要看它是否能处理这些真实分布问题，而不只是 few-shot benchmark 分数。
\end{warningbox}

\subsection{本章小结}

Few-shot image classification 给元学习提供了清晰实验场景：以 $N$-way $K$-shot episodes 训练学习算法，再测试其对新类别的快速适应。Omniglot 是典型 benchmark，但元学习思想也可迁移到语音、NLP 和其他任务中。

% =====================================================================
\section{总结与延伸}
% =====================================================================

P56 的核心不是介绍一个单独算法，而是把“学习算法里什么可以被学”系统展开。沿着梯度下降流程看，可学习对象从内到外不断扩大：
\begin{itemize}
    \item \textbf{初始化可以学：}MAML/Reptile 学一个适合快速适应的起点。
    \item \textbf{优化器可以学：}学习率、动量、更新规则和优化器状态都可以由经验决定。
    \item \textbf{结构可以学：}NAS 让 agent 或可微参数选择网络架构。
    \item \textbf{资料处理可以学：}data augmentation 与 sample weighting 不必完全靠人工试错。
    \item \textbf{分类流程可以学：}metric-based 方法把 support-query 比较本身变成可学习函数。
\end{itemize}

这也解释了“万物皆可 Meta”的准确含义：并不是所有问题都应该硬套元学习，而是当某个设计选择反复出现在许多任务中、且人工选择成本高时，就可以考虑把它提升为 task-level 学习对象。

\begin{importantbox}{判断一个组件是否值得 Meta 化}
    一个组件适合元学习，通常需要满足三个条件：它会显著影响任务表现；它在许多任务之间存在可复用规律；它的好坏可以通过 validation/query performance 反馈出来。若没有跨任务规律，meta-training 只会学到训练任务的偏好。
\end{importantbox}

从 P55 到 P56，可以形成一个完整检查表。看到任何元学习方法时，先问：
\begin{enumerate}
    \item 它学习的 $\phi$ 是什么？初始化、优化器、结构、资料处理、度量函数，还是其他？
    \item 它的 training tasks 与 testing tasks 是否分开？
    \item 它如何构造 support/query，如何计算 meta-loss？
    \item 它声称的优势来自 rapid adaptation、feature reuse、better optimization，还是更好的 data policy？
    \item 它是否真的泛化到新任务，还是只在同一 benchmark 上调得很好？
\end{enumerate}

这份检查表能帮助我们避免两个误区。第一，不要把 few-shot benchmark 成功自动等同于通用 learning to learn；第二，不要把所有预训练、多任务、NAS 或数据增强都粗略叫作元学习。真正关键的是跨任务训练结构与可迁移学习规则。

\subsection{本章小结}

P56 展示了元学习的广度：学习对象可以是初始化、优化器、架构、数据处理、样本权重或比较函数。它把 P55 的抽象框架具体化为一组方法族，也提醒我们：元学习的价值在于学习可迁移的“如何学习”，而不是给任何技巧贴上 meta 的标签。

\end{document}
